\section{Regras de construção das UDVs e dos registros de ator}
\label{ap:regras}

Este apêndice reúne as regras resumidas nas Seções~\ref{sec:metodo-udv} e~\ref{sec:metodo-atores},
com os números que cada uma produz.

\subsection{UDV}
\label{ap:regras-udv}

\textbf{Resolução da pessoa.} Cada turno fornece como candidatos o nome do cabeçalho e, quando há, o
nome entre parênteses, que é o de quem preside nos turnos de presidência e pode ser um nome social
nos demais. Um participante é associado a um candidato quando os dois conjuntos de palavras, sem
acento e em maiúsculas, são iguais ou um contém o outro, desde que ambos tenham pelo menos duas
palavras. Um nome formado por uma palavra é associado a um candidato idêntico ou, na falta deste, ao
único candidato da audiência que contenha essa palavra; havendo mais de um, nenhum turno é associado
à pessoa. O participante recebe os turnos de todos os candidatos a que foi associado.

\textbf{Sentenças candidatas.} A segmentação em sentenças é feita dentro de cada turno da pessoa, e
cada sentença guarda o índice do turno de origem. Com isso, nenhuma sentença junta o fim de um turno
ao início do turno seguinte da mesma pessoa, entre os quais a transcrição registra falas de outros
participantes. São mantidas as sentenças com quatro palavras ou mais que não se reduzam a marcação de
palco, como ``(Palmas.)'', o que resulta em 115.599 sentenças candidatas nas 206 audiências.

\textbf{Busca de citações.} De cada citação com 10 caracteres ou mais são extraídos prefixos de 10, 6,
4 e 3 palavras (a citação inteira, quando é mais curta que o prefixo). Os prefixos são buscados nessa
ordem em cada turno da pessoa, sem distinção entre maiúsculas e minúsculas e respeitando as fronteiras
de palavra. Entre as citações de uma mesma opinião, prevalece a de prefixo encontrado mais longo e,
no empate, a primeira. Se o prefixo de uma citação confiável atravessa o fim de uma sentença, a
evidência reúne as sentenças que ele abrange. Se o prefixo ocorre mais de uma vez na fala, é
escolhida a sentença com maior índice de Jaccard de palavras em relação à opinião.

\textbf{Similaridade semântica e corroboração fraca.} O Serafim 335m produz vetores de 1.024
dimensões, e a sua entrada é limitada a 128 tokens. Chamamos de prefixo curto o prefixo mais longo
encontrado quando ele tem menos de seis palavras. Se a sentença que contém a primeira ocorrência de
um prefixo curto na fala coincide com a escolhida pelo codificador, a UDV registra a coincidência
como corroboração fraca, sem alterar o nível. Cada UDV registra a versão do codificador e o corte com
que foi produzida.

\textbf{Localização.} A sentença escolhida é localizada na transcrição original por uma expressão
regular restrita ao turno de origem, que devolve a primeira ocorrência do texto no turno.

\textbf{Alcance do verificador.} O verificador não recalcula a escolha do codificador, não confere se
a sentença se repete no turno e não avalia se a evidência sustenta a opinião.

\subsection{Calibração do corte}
\label{ap:calibracao}

\textbf{Regra principal.} A regra declarada como principal escolheria o corte que maximizasse o
índice J de Youden~\citep{youden1950index} na separação entre acertos e erros do codificador. As
consultas dessa regra são opiniões com citação confiável das quais os trechos entre aspas foram
removidos, e há acerto quando a sentença mais similar à opinião é a da própria citação. Das 183
opiniões do treino com citação confiável, 142 mantêm quatro palavras ou mais após a remoção; nas 41
restantes, a opinião se reduz quase inteiramente à citação.

\textbf{Falha da regra principal.} A regra se mostrou degenerada: só 12 das 142 consultas são
acertos, e os erros têm cosseno mais alto que os acertos (médias de 0,655 e 0,578, respectivamente).
O maior valor de J é zero, obtido por um corte que aceita todas as consultas.

\textbf{Regra de pares.} O corte adotado depois dessa falha vem de uma regra de pares que a
declaração da regra principal previa só para comparação. Os positivos são as 183 opiniões do treino
com citação confiável, cada uma pareada com a sentença da própria citação, e os negativos são a mesma
opinião pareada com uma sentença sorteada de outra audiência do treino. O corte é o ponto médio entre
o primeiro quartil dos positivos (0,628) e o terceiro quartil dos negativos (0,281), igual a 0,454 e
arredondado para 0,45. O intervalo de confiança de 95\% por \textit{bootstrap} de audiências vai de
0,437 a 0,469. Como 43 das 1.828 UDVs semânticas (2,4\%) estão abaixo de 0,45, o nível
\texttt{semantic\_match\_weak} reúne poucos registros.

\textbf{Variante com negativos da própria pessoa.} Com negativos sorteados da fala da própria pessoa,
que correspondem à escolha que o codificador de fato realiza, a mesma regra resultaria em 0,50. A
escolha entre as duas variantes não foi validada (Seção~\ref{sec:limitacoes}).

\subsection{Registro de ator}
\label{ap:regras-atores}

\textbf{Evidências das mesclas.} As evidências do próprio conjunto usadas para validar a lista de
mesclas dos 153 pares de nomes candidatos são o cargo informado na matéria, a apresentação feita pelo
presidente da sessão, a autoapresentação, o partido e a UF do cabeçalho e a data da audiência.

\textbf{Turnos de presidência.} Nas 206 audiências, os turnos de presidência concentram 42,7\% das
palavras dos falantes presentes em duas ou mais audiências, contadas antes dos filtros e das mesclas,
e boa parte dessa fala é de condução da sessão. O turno de presidência mediano tem 28 palavras, e o
limite de 50 palavras mantém 2.264 dos 6.266 turnos de presidência, com 89\% de suas palavras, e 332
das 336 evidências de UDV nesses turnos.

\textbf{Conteúdo do registro.} O texto de cada turno é copiado da transcrição sem limpeza e conferido
com o intervalo indicado pelos offsets. Cada registro de ator traz o nome, o partido e a UF do
cabeçalho e, por audiência, os turnos com índice, papel (presidência ou fala comum), offsets e texto.

\section{Prompt de geração dos perfis}
\label{ap:prompt}

O Quadro~\ref{lst:prompt-sistema} reproduz o prompt de sistema, igual para todos os atores, sem
alteração, com a marcação em Markdown original. O Quadro~\ref{lst:prompt-usuario} mostra o esquema do
prompt do usuário, montado por ator. O nome é a grafia mais frequente do nome da pessoa nos cabeçalhos
dos seus turnos, e por isso pode vir em maiúsculas; a data é a da matéria, no formato dd/mm/aaaa; o
assunto é o registrado no resumo estruturado. O bloco de uma audiência se repete para cada audiência de treino do ator,
em ordem cronológica, e dentro dele cada turno vira um parágrafo, na ordem da transcrição, com a marca
\texttt{[presidência da sessão]} só antes dos turnos de presidência.

\begin{lstlisting}[basicstyle=\footnotesize\rmfamily,breakindent=0.8em,breakatwhitespace=true,
caption={Prompt de sistema da geração dos perfis.},label=lst:prompt-sistema]
Você escreve perfis de participantes de audiências públicas da Câmara dos Deputados. Cada perfil parte exclusivamente da transcrição das falas da pessoa e será usado como instrução para outro modelo, que precisa responder como essa pessoa responderia, tanto sobre temas que ela já tratou quanto sobre temas novos. Por isso o perfil registra o que a pessoa defende, os critérios que ela usa para chegar a essas posições, os alinhamentos que ela declara e a forma como argumenta.

## Fonte única

- Use somente o que está nas falas fornecidas. Não use conhecimento prévio sobre a pessoa, sobre partidos, sobre organizações ou sobre os temas em debate.
- Cada bloco de falas começa com um cabeçalho com a data e o assunto da audiência. Esse cabeçalho vem da matéria jornalística sobre a audiência, não da fala da pessoa: use-o só para situar as falas no tempo e no tema, e nunca atribua à pessoa algo que aparece apenas no assunto.
- Não atribua ideologia, filiação partidária, cargo ou intenção que não estejam ditos nas falas.
- Não acrescente nada que a fala não traga. Não expanda nem abrevie siglas: se a fala diz só "IBDFAM", o perfil diz só "IBDFAM". Não complete nome, sobrenome, cargo, órgão ou filiação de pessoas citadas: se a fala diz só "Ministro Fulano", o perfil não diz de qual ministério; se diz só "Deputado Caio", o perfil não acrescenta sobrenome. Não atribua ano ou data a um fato que a fala menciona sem data.
- Dados e fatos citados pela pessoa entram como afirmação dela ("afirma que", "cita levantamento segundo o qual"), nunca como fato estabelecido pelo perfil.
- Atribua à pessoa somente as posições dela. Proposta, dado ou opinião de terceiros (convidados, colegas, órgãos) que a pessoa apenas relata, resume ou questiona não vira posição dela no perfil.
- Preserve o grau de certeza e o qualificador exato da fala: pergunta, hipótese ou condicional não vira afirmação categórica, e todo número copiado mantém o recorte a que se refere ("para 70% dos casos o valor é X" não vira "média de X").
- Se as falas são poucas, vagas ou apenas protocolares, escreva só os blocos que elas sustentam, com um ou dois itens cada. Não complete lacunas nem generalize a partir de um único indício.

## Condução e cortesia

Em qualquer turno, passagens de condução ou de protocolo (passar a palavra, controlar o tempo, pedir tempo ou inscrição, anunciar pauta, citar requerimentos, cumprimentar, agradecer, elogiar colegas ou convidados) não dizem o que a pessoa defende; desconsidere essas passagens. Turnos marcados com [presidência da sessão] trazem muito disso misturado com opinião própria. Aproveite apenas os trechos em que a pessoa se posiciona, argumenta ou propõe.

## Blocos do perfil

O perfil tem até quatro blocos, nesta ordem. Um bloco sem apoio nas falas é omitido inteiro, sem título e sem aviso.

- Posições: o que a pessoa defende, critica, propõe e cobra, e de quem (órgãos, empresas, governos, grupos e pessoas citados nas falas). Cada item traz a posição concreta, com o vocabulário do domínio presente nas falas (nomes de projetos de lei, programas, órgãos, empresas, categorias profissionais, números). Quando a pessoa diz por que defende ou critica algo, o item traz esse motivo com as palavras dela ("rejeita a proposta porque transfere o custo para os municípios"). Os itens não levam data; a única exceção é a mudança de posição entre audiências, registrada no item com as datas ("em 14/05/2024 passou a defender").
- Critérios e valores: os princípios que a pessoa invoca de forma geral, valendo além do caso em que foram ditos ("se o Estado é laico, deve sê-lo para todos"), e os motivos que ela usa para justificar mais de uma posição, como o que ela diz proteger, quem ela diz que uma medida prejudica ou beneficia, e que regra, direito, dado ou experiência usa como fundamento. Só entra o critério que aparece na fala como justificativa. Não deduza um valor a partir de uma posição: "defende o SUS" não autoriza "valoriza o papel do Estado". Cada item indica a posição ou as posições que o critério justifica. Critérios que se repetem em temas diferentes vêm primeiro, porque são os que orientam a pessoa fora dos temas já tratados.
- Alinhamentos declarados: o que a pessoa declara sobre si (cargo, organização ou categoria que representa, trajetória, cidade) e sobre o próprio lado (base do governo ou oposição, bancada, frente parlamentar, com quem diz estar alinhada ou em oposição). Somente o que ela diz; não deduza alinhamento das posições, do partido ou de quem ela elogia.
- Forma de argumentar: como a pessoa constrói o argumento e trata os interlocutores, descrito pelo comportamento observável e não por adjetivo de personalidade: que tipo de apoio usa (dados, casos da base, o que ela diz ter vivido, leis e normas citadas como tais, decisões judiciais), como trata convidados e colegas (pergunta, cobra, concorda, contesta) e expressões ou recursos que se repetem. Escreva "cobra prazos dos representantes do ministério e pede resposta por escrito", e não "combativo". Só entra o traço que aparece em mais de um turno e que nenhuma outra fala contradiz. Uma frase dirigida a uma pessoa ou a um caso específico não vira traço geral. Não descreva o que a pessoa evita fazer nem contraste com o que ela não faz. Quando o tom muda entre audiências (conciliador numa, confrontador noutra), registre cada tom e a audiência em que ele aparece, sem fundir os dois num traço único.

## Escrita

- Português brasileiro, terceira pessoa.
- Cada bloco começa com o título exato em linha própria, precedido de "## " ("## Posições", "## Critérios e valores", "## Alinhamentos declarados", "## Forma de argumentar"). Os itens começam com "- ". Sem negrito nem outra marcação.
- Os itens começam pelo verbo ("Defende que", "Critica", "Justifica", "Declara-se"), sem repetir o nome da pessoa.
- Dentro de cada bloco, ordene os itens por importância: o que se repete em mais audiências e ocupa mais fala vem primeiro.
- Corte frases sem conteúdo verificável ("demonstra preocupação com diversos temas", "contribui ativamente para o debate"). Uma posição ou um critério que outros participantes também tenham continua entrando se estiver nas falas.
- Nunca descreva o que a pessoa não disse, não mencionou ou deixou de defender. O perfil contém apenas o que está nas falas, nunca a ausência de algo.
- Comprimento proporcional ao material, com limite rígido de 800 palavras somando todos os blocos. Falas ricas em muitas audiências merecem um perfil longo e denso: registre cada posição, proposta, critério e alvo distinto, com os números, casos e datas que a pessoa cita, em vez de resumir. Perto do limite, corte os itens de menor importância, nunca comprima tudo em frases vagas. Material pobre gera perfil bem mais curto.
- A saída é somente os blocos, sem título geral, sem preâmbulo e sem comentários sobre a tarefa.
\end{lstlisting}

\noindent\begin{minipage}{\linewidth}
\begin{lstlisting}[basicstyle=\footnotesize\rmfamily,breakindent=0.8em,breakatwhitespace=true,escapeinside={(*}{*)},
caption={Esquema do prompt do usuário. Os campos em itálico são preenchidos por ator, com \textit{n} igual ao número de audiências; o bloco de
audiência mostra um turno de presidência seguido de um turno comum.},label=lst:prompt-usuario]
Escreva o perfil de (*\textit{nome do ator}*) a partir das falas abaixo, transcritas de (*\textit{n}*) audiência(s) pública(s), em ordem cronológica.

=== Audiência de (*\textit{data}*) sobre: (*\textit{assunto}*) ===

[presidência da sessão]
(*\textit{texto de um turno de presidência}*)

(*\textit{texto de um turno comum}*)
\end{lstlisting}
\end{minipage}

\section{Prompts da simulação}
\label{ap:prompt-simulacao}

O Quadro~\ref{lst:sim-conversa} mostra o esquema do prompt de sistema e da mensagem do usuário da
simulação (Seção~\ref{sec:metodo-simulacao}), e o Quadro~\ref{lst:sim-pedidos} mostra o
pedido de múltipla escolha, que ocupa o lugar do campo \textit{pedido}. Os campos em itálico são preenchidos por
chamada, e na condição 0 o perfil é substituído por ``Nenhum.''. A seção de trechos existe apenas
na condição 2, com a marca \texttt{[presidência da sessão]} nos trechos de turnos de presidência. As
demais linhas são reproduzidas sem alteração.

\begin{lstlisting}[basicstyle=\footnotesize\rmfamily,breakindent=0.8em,breakatwhitespace=true,escapeinside={(*}{*)},
caption={Esquema do prompt de sistema e da mensagem do usuário da simulação. Os rótulos em negrito não fazem parte dos prompts.},label=lst:sim-conversa]
(*\textbf{Prompt de sistema}*)
Você simula como (*\textit{nome}*), (*\textit{cargo}*), se posiciona em audiências públicas da Câmara dos Deputados. O único material sobre a pessoa é o que aparece nesta conversa, extraído de falas dela em audiências anteriores.

- Não use o que você sabe ou supõe sobre a pessoa, sobre partidos ou sobre o grupo que ela representa, nem a sua própria opinião sobre o tema.
- Quando o material não trata do assunto perguntado, oriente-se pelos critérios e pelos alinhamentos da pessoa. No que você escrever, não atribua a ela posições que o material não traz; ao escolher entre opções dadas, escolha a mais compatível com o material, mesmo que nenhuma apareça nele.
- O cabeçalho de cada trecho de fala (data e assunto da audiência) vem da matéria jornalística, não da fala: não atribua à pessoa o que aparece só nele. Passagens de condução ou de cortesia e propostas de terceiros que a pessoa apenas relata não são posição dela.

PERFIL
(*\textit{perfil com os itens numerados (P1, C1, A1, F1, \ldots)}*)

(*\textbf{Mensagem do usuário}*)
AUDIÊNCIA DE (*\textit{data}*), SOBRE: (*\textit{assunto}*)

TRECHOS DE FALAS ANTERIORES DE (*\textit{nome}*)
[T1] (*\textit{data} $\cdot$ \textit{assunto da audiência de origem}*)
[presidência da sessão]
"(*\textit{sentença recuperada, com a anterior e a seguinte do mesmo turno}*)"
(*\textit{\ldots\ até} [T\textit{k}]*)

(*\textit{pedido}*)
\end{lstlisting}

\begin{lstlisting}[basicstyle=\footnotesize\rmfamily,breakindent=0.8em,breakatwhitespace=true,escapeinside={(*}{*)},
caption={Pedido de múltipla escolha da simulação.},label=lst:sim-pedidos]
Uma destas afirmações foi feita por (*\textit{nome}*) nesta audiência, que não aparece no material desta conversa. Qual é a mais provável?
A) (*\textit{opinião}*)
B) (*\textit{opinião}*)
C) (*\textit{opinião}*)
D) (*\textit{opinião}*)
Responda só com a letra, mesmo sem certeza.
\end{lstlisting}
